% ---- IBD0 ----
% Code for the TikZ picture showing how the Young et al (2022) imputation procedure works.
\begin{subfigure}[b]{0.25\textwidth}
    \centering
    \begin{tikzpicture}
        % Draw parental nodes
        \node[state,dashed,gray] (mother) at (-1, 3) {?};
        \node[state,dashed,gray] (father) at (1, 3)  {?};
        \node [above=0.25cm of mother] {\textcolor{orange}{Mother}};
        \node [above=0.25cm of father] {\textcolor{blue}{Father}};
        % Draw individual, and their sibling, nodes.
        \node[state] (individual) [below=1cm of mother] {1};
        \node[state] (sibling)    [below=1cm of father] {1};
        \node [below=0.25cm of individual] {\textcolor{ForestGreen}{Child}};
        \node [below=0.25cm of sibling]    {\textcolor{purple}{Sibling}};
        % Draw the family connections
        \coordinate (betweenparents) at ($(mother)!0.5!(father)$);
        \coordinate (aboveindividual) at ($(mother)!0.5!(individual)$);
        \coordinate (abovesibling) at ($(father)!0.5!(sibling)$);
        \coordinate (aboveboth) at ($(aboveindividual)!0.5!(abovesibling)$);
        \path[-,dashed,gray] (mother) edge (father);
        \path[-,dashed,gray] (betweenparents) edge (aboveboth);
        \path[-] (aboveindividual) edge (abovesibling);
        \path[-] (aboveindividual) edge (individual);
        \path[-] (abovesibling) edge (sibling);
        % Bridge to the others
        \node[scale=2] [right=1cm of abovesibling]{$\Rightarrow$};
        \node [left=3.5cm of abovesibling]{IB0:};
    \end{tikzpicture}
\end{subfigure}
\hskip3cm
\begin{subfigure}[b]{0.25\textwidth}
    \centering
    \begin{tikzpicture}
        % Draw parental nodes
        \node[state] (mother) at (-1, 3) {1};
        \node[state] (father) at (1, 3)  {1};
        \node [above=0.25cm of mother] {\textcolor{orange}{Mother}};
        \node [above=0.25cm of father] {\textcolor{blue}{Father}};
        \node [color=white,below=0.25cm of individual] {Child};
        \node [color=white,below=0.25cm of sibling]    {Sibling};
        % Draw individual, and their sibling, nodes.
        \node[color=white,state] (individual) [below=1cm of mother] {1};
        \node[color=white,state] (sibling)    [below=1cm of father] {0};
        % Draw the family connections
        \coordinate (betweenparents) at ($(mother)!0.5!(father)$);
        \coordinate (aboveindividual) at ($(mother)!0.5!(individual)$);
        \coordinate (abovesibling) at ($(father)!0.5!(sibling)$);
        \coordinate (aboveboth) at ($(aboveindividual)!0.5!(abovesibling)$);
        \path[-,dashed,gray] (mother) edge (father);
        \path[-,dashed,gray] (betweenparents) edge (aboveboth);
        %\path[-] (aboveindividual) edge (abovesibling);
        %\path[-] (aboveindividual) edge (individual);
        %\path[-] (abovesibling) edge (sibling);
        \node [below=0.5cm of aboveboth] {Mean parent SNP $= 1$};
    \end{tikzpicture}
\end{subfigure}

% ---- IBD1 ----
% Code for the TikZ picture showing how the Young et al (2022) imputation procedure works.
\begin{subfigure}[b]{0.25\textwidth}
    \centering
    \begin{tikzpicture}
        % Draw parental nodes
        \node[state,dashed,gray] (mother) at (-1, 3) {?};
        \node[state,dashed,gray] (father) at (1, 3)  {?};
        \node [above=0.25cm of mother] {\textcolor{orange}{Mother}};
        \node [above=0.25cm of father] {\textcolor{blue}{Father}};
        % Draw individual, and their sibling, nodes.
        \node[state] (individual) [below=1cm of mother] {2};
        \node[state] (sibling)    [below=1cm of father] {1};
        \node [below=0.25cm of individual] {\textcolor{ForestGreen}{Child}};
        \node [below=0.25cm of sibling]    {\textcolor{purple}{Sibling}};
        % Draw the family connections
        \coordinate (betweenparents) at ($(mother)!0.5!(father)$);
        \coordinate (aboveindividual) at ($(mother)!0.5!(individual)$);
        \coordinate (abovesibling) at ($(father)!0.5!(sibling)$);
        \coordinate (aboveboth) at ($(aboveindividual)!0.5!(abovesibling)$);
        \path[-,dashed,gray] (mother) edge (father);
        \path[-,dashed,gray] (betweenparents) edge (aboveboth);
        \path[-] (aboveindividual) edge (abovesibling);
        \path[-] (aboveindividual) edge (individual);
        \path[-] (abovesibling) edge (sibling);
        % Bridge to the others
        \node[scale=2] [right=1cm of abovesibling]{$\Rightarrow$};
        \node [left=3.5cm of abovesibling]{IB1:};
    \end{tikzpicture}
\end{subfigure}
\hskip3cm
\begin{subfigure}[b]{0.25\textwidth}
    \centering
    \begin{tikzpicture}
        % Draw parental nodes
        \node[state] (mother) at (-1, 3) {1};
        \node[state,dashed,gray] (father) at (1, 3)  {?};
        \node [above=0.25cm of mother] {\textcolor{orange}{Mother}};
        \node [above=0.25cm of father] {\textcolor{blue}{Father}};
        \node [color=white,below=0.25cm of individual] {Child};
        \node [color=white,below=0.25cm of sibling]    {Sibling};
        % Draw individual, and their sibling, nodes.
        \node[color=white,state] (individual) [below=1cm of mother] {2};
        \node[color=white,state] (sibling)    [below=1cm of father] {1};
        % Draw the family connections
        \coordinate (betweenparents) at ($(mother)!0.5!(father)$);
        \coordinate (aboveindividual) at ($(mother)!0.5!(individual)$);
        \coordinate (abovesibling) at ($(father)!0.5!(sibling)$);
        \coordinate (aboveboth) at ($(aboveindividual)!0.5!(abovesibling)$);
        \path[-,dashed,gray] (mother) edge (father);
        \path[-,dashed,gray] (betweenparents) edge (aboveboth);
        %\path[-] (aboveindividual) edge (abovesibling);
        %\path[-] (aboveindividual) edge (individual);
        %\path[-] (abovesibling) edge (sibling);
        \node [right=0.4cm of father] {or};
        \node [below=0.5cm of aboveboth] {Mean = $\frac{1 + f}{2}$};
    \end{tikzpicture}
\end{subfigure}
\begin{subfigure}[b]{0.2\textwidth}
    \centering
    \begin{tikzpicture}
        % Draw parental nodes
        \node[state] (mother) at (-1, 3) {2};
        \node[state,dashed,gray] (father) at (1, 3)  {?};
        \node [above=0.25cm of mother] {\textcolor{orange}{Mother}};
        \node [above=0.25cm of father] {\textcolor{blue}{Father}};
        \node [color=white,below=0.25cm of individual] {Child};
        \node [color=white,below=0.25cm of sibling]    {Sibling};
        % Draw individual, and their sibling, nodes.
        \node[color=white,state] (individual) [below=1cm of mother] {1};
        \node[color=white,state] (sibling)    [below=1cm of father] {0};
        % Draw the family connections
        \coordinate (betweenparents) at ($(mother)!0.5!(father)$);
        \coordinate (aboveindividual) at ($(mother)!0.5!(individual)$);
        \coordinate (abovesibling) at ($(father)!0.5!(sibling)$);
        \coordinate (aboveboth) at ($(aboveindividual)!0.5!(abovesibling)$);
        \path[-,dashed,gray] (mother) edge (father);
        \path[-,dashed,gray] (betweenparents) edge (aboveboth);
        %\path[-] (aboveindividual) edge (abovesibling);
        %\path[-] (aboveindividual) edge (individual);
        %\path[-] (abovesibling) edge (sibling);
        \node [below=0.5cm of aboveboth] {Mean = $\frac{2 + f}{2}$};
    \end{tikzpicture}
\end{subfigure}

% ---- IBD0 ----
% Code for the TikZ picture showing how the Young et al (2022) imputation procedure works.
\begin{subfigure}[b]{0.25\textwidth}
    \centering
    \begin{tikzpicture}
        % Draw parental nodes
        \node[state,dashed,gray] (mother) at (-1, 3) {?};
        \node[state,dashed,gray] (father) at (1, 3)  {?};
        \node [above=0.25cm of mother] {\textcolor{orange}{Mother}};
        \node [above=0.25cm of father] {\textcolor{blue}{Father}};
        % Draw individual, and their sibling, nodes.
        \node[state] (individual) [below=1cm of mother] {1};
        \node[state] (sibling)    [below=1cm of father] {1};
        \node [below=0.25cm of individual] {\textcolor{ForestGreen}{Child}};
        \node [below=0.25cm of sibling]    {\textcolor{purple}{Sibling}};
        % Draw the family connections
        \coordinate (betweenparents) at ($(mother)!0.5!(father)$);
        \coordinate (aboveindividual) at ($(mother)!0.5!(individual)$);
        \coordinate (abovesibling) at ($(father)!0.5!(sibling)$);
        \coordinate (aboveboth) at ($(aboveindividual)!0.5!(abovesibling)$);
        \path[-,dashed,gray] (mother) edge (father);
        \path[-,dashed,gray] (betweenparents) edge (aboveboth);
        \path[-] (aboveindividual) edge (abovesibling);
        \path[-] (aboveindividual) edge (individual);
        \path[-] (abovesibling) edge (sibling);
        % Bridge to the others
        \node[scale=2] [right=1cm of abovesibling]{$\Rightarrow$};
        \node [left=3.5cm of abovesibling]{IB0:};
    \end{tikzpicture}
\end{subfigure}
\hskip3cm
\begin{subfigure}[b]{0.25\textwidth}
    \centering
    \begin{tikzpicture}
        % Draw parental nodes
        \node[state,dashed,gray] (mother) at (-1, 3) {?};
        \node[state,dashed,gray] (father) at (1, 3)  {?};
        \node [above=0.25cm of mother] {\textcolor{orange}{Mother}};
        \node [above=0.25cm of father] {\textcolor{blue}{Father}};
        \node [color=white,below=0.25cm of individual] {Child};
        \node [color=white,below=0.25cm of sibling]    {Sibling};
        % Draw individual, and their sibling, nodes.
        \node[color=white,state] (individual) [below=1cm of mother] {1};
        \node[color=white,state] (sibling)    [below=1cm of father] {0};
        % Draw the family connections
        \coordinate (betweenparents) at ($(mother)!0.5!(father)$);
        \coordinate (aboveindividual) at ($(mother)!0.5!(individual)$);
        \coordinate (abovesibling) at ($(father)!0.5!(sibling)$);
        \coordinate (aboveboth) at ($(aboveindividual)!0.5!(abovesibling)$);
        \path[-,dashed,gray] (mother) edge (father);
        \path[-,dashed,gray] (betweenparents) edge (aboveboth);
        %\path[-] (aboveindividual) edge (abovesibling);
        %\path[-] (aboveindividual) edge (individual);
        %\path[-] (abovesibling) edge (sibling);
        \node [below=0.5cm of aboveboth] {Mean parent SNP $= \frac{1 + 2 f}{2}$};
    \end{tikzpicture}
\end{subfigure}
